\documentclass[10pt,landscape]{article}

\usepackage[margin=0.32in]{geometry}
\usepackage[T1]{fontenc}
\usepackage[utf8]{inputenc}
\usepackage{lmodern}
\usepackage{textcomp}
\usepackage{microtype}
\usepackage[table,dvipsnames]{xcolor}
\usepackage{array}
\usepackage{booktabs}
\usepackage{enumitem}
\usepackage{hyperref}
\usepackage{listings}
\usepackage{multicol}
\usepackage{tabularx}
\usepackage{titlesec}

\pagestyle{empty}
\setlength{\parindent}{0pt}
\setlength{\parskip}{1.5pt}
\setlength{\columnsep}{13pt}
\setlength{\columnseprule}{0.2pt}
\setcounter{secnumdepth}{0}

\definecolor{ToshBlue}{HTML}{0F4C81}
\definecolor{ToshInk}{HTML}{1F2937}
\definecolor{ToshMuted}{HTML}{5B6470}
\definecolor{ToshPanel}{HTML}{F4F7FB}
\definecolor{ToshRule}{HTML}{CBD5E1}
\definecolor{ToshCodeBg}{HTML}{F8FAFC}

\color{ToshInk}

\titleformat{\section}
  {\large\bfseries\color{ToshBlue}}
  {}
  {0pt}
  {}
\titlespacing*{\section}{0pt}{0.4em}{0.12em}

\setlist[itemize]{leftmargin=1.05em,itemsep=0.5pt,topsep=1.5pt,parsep=0pt,partopsep=0pt}

\lstdefinestyle{tosh}{
  basicstyle=\ttfamily\scriptsize,
  backgroundcolor=\color{ToshCodeBg},
  frame=single,
  rulecolor=\color{ToshRule},
  framerule=0.4pt,
  xleftmargin=0pt,
  xrightmargin=0pt,
  aboveskip=0.18em,
  belowskip=0.24em,
  columns=fullflexible,
  keepspaces=true,
  showstringspaces=false,
  breaklines=true,
  literate={°}{{\textdegree}}1
           {μ}{{\ensuremath{\mu}}}1
           {Ω}{{\ensuremath{\Omega}}}1
}
\lstset{style=tosh}

\newcommand{\sheettitle}[3]{
  {\LARGE\bfseries\color{ToshBlue} ToSh Cheat Sheet: #1\par}
  \vspace{0.08em}
  {\small\color{ToshMuted} #2 \hfill #3\par}
  \vspace{0.18em}
  {\color{ToshRule}\rule{\linewidth}{0.8pt}}
  \vspace{0.12em}
}

\newcommand{\note}[1]{
  \vspace{0.12em}
  \fcolorbox{ToshRule}{ToshPanel}{%
    \parbox{\dimexpr\linewidth-2\fboxsep-2\fboxrule\relax}{\footnotesize #1}%
  }
  \vspace{0.12em}
}

\newcommand{\tighttable}[2]{
  {\footnotesize
  \renewcommand{\arraystretch}{1.08}
  \begin{tabularx}{\linewidth}{@{}#1@{}}
  #2
  \end{tabularx}}
}


\begin{document}
\footnotesize
\sheettitle{Control Flow And Matching}{Branching, loops, errors, cleanup, and expression-style matching}{parser/tests first}

\begin{multicols*}{3}

\section{If And Conditional Style}
\begin{lstlisting}
if ($age >= 18) {
    echo adult
} else if ($age >= 13) {
    echo teen
} else {
    echo child
}
\end{lstlisting}

\begin{lstlisting}
var label = if ($count == 1) { "item" } else { "items" }
echo $"You have {$count} {$label}."
\end{lstlisting}

\note{\texttt{if} works both as a statement and as an expression, which makes it handy inside assignments and interpolation.}

\section{Loops}
\begin{lstlisting}
for item in (ls | first 3) {
    echo $item.Name
}

for i in 1..5 {
    echo $i
}
\end{lstlisting}

\begin{lstlisting}
while (($count < 3)) {
    echo $count
    $count += 1
}

until (($ready)) {
    sleep 1s
}
\end{lstlisting}

\tighttable{>{\ttfamily\raggedright\arraybackslash}p{0.30\linewidth}X}{
\toprule
Keyword & Use \\
\midrule
\texttt{break} & Exit the current loop. \\
\texttt{continue} & Skip to the next iteration. \\
\texttt{return} & Exit the current function/script. \\
\texttt{throw} & Raise an error value. \\
\bottomrule
}

\columnbreak

\section{Errors And Cleanup}
\begin{lstlisting}
try {
    read-file config.json | from json
} catch (err) {
    writeline $"Error: {$err.Message}"
} finally {
    writeline done
}
\end{lstlisting}

\begin{lstlisting}
func process-file(path) {
    var handle = (open-file --read $path)
    defer { close $handle }

    return (read-to-end $handle)
}
\end{lstlisting}

\begin{itemize}
\item \texttt{catch (err)} binds the thrown value or exception object.
\item \texttt{finally} always runs.
\item \texttt{defer} schedules cleanup when the current scope exits.
\end{itemize}

\section{Switch}
\begin{lstlisting}
switch ($state) {
    case file { echo file }
    case dir  { echo dir }
    default   { echo other }
}
\end{lstlisting}

\begin{lstlisting}
switch ($temperature) {
    case < 32   { echo Freezing }
    case 32..50 { echo Cold }
    case > 75   { echo Hot }
}
\end{lstlisting}

\note{\texttt{switch} is statement-style and does not fall through. Each \texttt{case} is independent; no extra \texttt{break} is needed.}

\section{Match Expressions}
\begin{lstlisting}
var label = match ($status) {
    200 => "OK"
    404 => "Not Found"
    default => "Other"
}
\end{lstlisting}

\begin{lstlisting}
match ($value) {
    _ is int => "integer"
    100..199 => "informational"
    _ =~ "\\.csv$" => "csv"
    default => "fallback"
}
\end{lstlisting}

\columnbreak

\section{Guards And Ordering}
\begin{lstlisting}
match ($n) {
    0 => "zero"
    _ > 0 if (($n % 2 == 0)) => "positive even"
    _ > 0 => "positive odd"
    default => "other"
}
\end{lstlisting}

\begin{itemize}
\item Arms are tested top to bottom.
\item The first matching arm wins.
\item \texttt{default} is required for exhaustive behavior.
\item Guards use \texttt{if (...)} after the pattern.
\end{itemize}

\section{Control Patterns In Scripts}
\begin{lstlisting}
func recent(span: TimeSpan) =>
    ls -la | where _.Modified > ((date now) - $span)

func big-files(min: StorageSize) {
    ls -la | where _.Type == file | where _.Size >= $min
}
\end{lstlisting}

\begin{lstlisting}
echo $"2 + 2 = {2 + 2}"
echo $"Files: {ls | where _.Type == file | count}"
echo ((date now) - 2d)
\end{lstlisting}

\section{Helpful Mental Model}
\begin{itemize}
\item Use \texttt{if} when you want branch-oriented control flow.
\item Use \texttt{match} when you want a value back from a compact pattern table.
\item Use \texttt{switch} when each arm performs statements or side effects.
\item Pair loops with pipeline commands whenever you can; fall back to loops when stateful control is clearer.
\end{itemize}

\note{ToSh scripts stay readable when most data reshaping happens in pipelines, while loops and \texttt{match} handle the control decisions around those pipelines.}

\end{multicols*}
\newpage

\sheettitle{Control Flow And Matching}{Pattern catalog, defer ordering, and practical control recipes}{page 2 / reference}

\begin{multicols*}{3}

\section{Pattern Catalog}
\begin{lstlisting}
200 => "OK"                  # exact value
_ is int => "integer"        # runtime type
_ >= 90 => "A"               # comparison
100..199 => "info"           # range
_ =~ "\\.csv$" => "csv"      # regex
_ > 0 if (($n % 2 == 0)) => "even"
\end{lstlisting}

\begin{itemize}
\item Match arms are checked top to bottom.
\item The first matching arm wins.
\item Keep \texttt{default} last and obvious.
\end{itemize}

\section{Match Recipes}
\begin{lstlisting}
func http_category(status: int) -> string {
    return match ($status) {
        100..199 => "Informational"
        200..299 => "Success"
        300..399 => "Redirection"
        400..499 => "Client Error"
        500..599 => "Server Error"
        default  => "Unknown"
    }
}
\end{lstlisting}

\begin{lstlisting}
match ($file) {
    _ =~ "\\.pdf$" => "document"
    _ =~ "\\.(csv|tsv)$" => "data"
    _ =~ "\\.tar\\." => "archive"
    default => "other"
}
\end{lstlisting}

\section{Switch Recipes}
\begin{lstlisting}
switch ($state) {
    case Inv.StockState.Low { echo low }
    case Inv.StockState.Out { echo out }
    default { echo ok }
}
\end{lstlisting}

\columnbreak

\section{Try, Throw, Return}
\begin{lstlisting}
try {
    throw "boom"
} catch (err) {
    echo $err
} finally {
    writeline done
}
\end{lstlisting}

\begin{lstlisting}
func find-first(items, predicate) {
    for item in $items {
        if (invoke $predicate $item) {
            return $item
        }
    }
    return null
}
\end{lstlisting}

\begin{itemize}
\item \texttt{return} can bubble out from nested loops and blocks.
\item Unhandled \texttt{throw} becomes a top-level shell diagnostic.
\item \texttt{catch (err)} gives you the thrown value or exception object.
\end{itemize}

\section{Defer Patterns}
\begin{lstlisting}
func process-file(path) {
    var handle = (open-file --read $path)
    defer { close $handle }

    return (read-to-end $handle)
}
\end{lstlisting}

\begin{lstlisting}
defer { echo "third" }
defer { echo "second" }
defer { echo "first" }
\end{lstlisting}

\note{\texttt{defer} runs in LIFO order. It is the easiest way to clean up file handles, temp resources, or partial state when a scope exits early.}

\section{Loop Recipes}
\begin{lstlisting}
for proc in (ps | sort Memory | reverse | first 3) {
    writeline $proc.Name
}
\end{lstlisting}

\begin{lstlisting}
echo one skip two three | each {
    if ((_ == "skip")) { continue }
    echo _
}
\end{lstlisting}

\columnbreak

\section{Real Script Patterns}
\begin{lstlisting}
ls -la | first 10 | each {
    var tag = match (_.Type.ToString()) {
        "dir" => "DIR "
        "file" if ((_.Size > 1kb)) => "BIG "
        "file" => "file"
        "symlink" => "link"
        default => "????"
    }
    echo $"[{$tag}] {_.Name}"
}
\end{lstlisting}

\begin{lstlisting}
func classify_score(score: int) -> string {
    return match ($score) {
        _ >= 90 => "A"
        _ >= 80 => "B"
        _ >= 70 => "C"
        _ >= 60 => "D"
        default => "F"
    }
}
\end{lstlisting}

\section{Choose The Form}
\tighttable{>{\ttfamily\raggedright\arraybackslash}p{0.24\linewidth}X}{
\toprule
Use & When it fits best \\
\midrule
\texttt{if} & one-off branching with side effects or short expressions \\
\texttt{match} & compact value selection from ordered patterns \\
\texttt{switch} & statement blocks keyed from one tested value \\
\texttt{loops} & stateful iteration when pipelines are not enough \\
\texttt{try}/\texttt{defer} & failure handling and cleanup \\
\bottomrule
}

\note{A readable ToSh script usually keeps data reshaping in pipelines, then uses \texttt{match}, \texttt{if}, and loops only where control flow really matters.}

\end{multicols*}
\end{document}
